\documentclass{article}
\usepackage[utf8]{inputenc}
\usepackage{mathtools,amssymb}
\usepackage{lipsum}
\usepackage{hyperref}
\usepackage{layout}
\usepackage{geometry}
\usepackage{listings}
\usepackage{color}
\usepackage{xcolor}
\usepackage{amsthm}
\usepackage{booktabs}
\usepackage{array}
\newcommand*{\vertbar}[1][10ex]{\rule[-1ex]{0.5pt}{#1}}
\newcommand*{\horzbar}{\rule[.5ex]{2.5ex}{0.5pt}}
\newcommand{\norm}[2]{\left\lVert#1\right\rVert_#2}
\newcommand{\abs}[1]{\lvert#1\rvert}

\title{6.S974 class notes}
\author{Erich Trieschman}
\date{2026 Fall semester}

% compact spacing
\usepackage{setspace}\setstretch{0.95}
\usepackage{parskip}\setlength{\parskip}{4pt}
\setlength{\parindent}{0pt}
\usepackage{enumitem}
\setlist[itemize]{topsep=2pt, itemsep=2pt, parsep=0pt, partopsep=0pt}
\setlist[enumerate]{topsep=2pt, itemsep=2pt, parsep=0pt, partopsep=0pt}
\usepackage{titlesec}
\titlespacing*{\section}{0pt}{3.5ex plus 0.3ex minus 0.3ex}{0.5ex}
\titlespacing*{\subsection}{0pt}{2.5ex plus 0.2ex minus 0.2ex}{0.3ex}
\titlespacing*{\subsubsection}{0pt}{1.5ex plus 0.2ex minus 0.2ex}{0.2ex}

\newcommand{\userMarginInMm}{15}
\geometry{
%  legal,
 left=\userMarginInMm mm,
 right=\userMarginInMm mm,
 top=\userMarginInMm mm,
 bottom=\userMarginInMm mm,
 footskip=5mm}


 \hypersetup{
    colorlinks=true,
    linkcolor=blue,
    filecolor=magenta,
    urlcolor=cyan,
    pdfpagemode=FullScreen,
}

\definecolor{dkgreen}{rgb}{0,0.6,0}
\definecolor{gray}{rgb}{0.5,0.5,0.5}
\definecolor{mauve}{rgb}{0.58,0,0.82}
\newcommand{\todo}[1]{\textit{\textcolor{orange}{ [#1]}}}

% \new marks a correction to a statement in the first version of these notes.
% To remove the color: \renewcommand{\new}[1]{#1}
\newcommand{\new}[1]{\textcolor{blue}{#1}}
% \res labels a named result.
\newcommand{\res}[1]{\textbf{#1.}}

\lstset{frame=tb, language=Python,
    aboveskip=3mm,
  belowskip=3mm,
  showstringspaces=false,
  columns=flexible,
  basicstyle={\small\ttfamily},
  numbers=none,
  numberstyle=\tiny\color{gray},
  keywordstyle=\color{blue},
  commentstyle=\color{dkgreen},
  stringstyle=\color{mauve},
  breaklines=true,
  breakatwhitespace=true,
  tabsize=3}

\begin{document}
\maketitle
\tableofcontents

\section*{Exam guidance}
The exam will cover the material from Lectures 2-20 and Problem Sets 1-3 except:
\begin{itemize}
    \item Slides 35-36 of Lectures 7-8 (Intuitive criterion)
    \item Proofs of theorems in Lecture 12 (Learning in games)
    \item Slides 52-60 of Lectures 16-17 (Stopping rule and Rothschild's theorem)
    \item Slides 23-57 of Lecture 18 (Markov games)
\end{itemize}
The exam will contain 4-5 problems, which will be much shorter than in the problem sets. They will test whether you understand the concepts and can apply them in a new setting. There will be no coding question. If you understand the problem sets and the lecture examples (and you can solve them again on your own), then you are well prepared!

Cheatsheet: 1 page, two-sided.

%======================================================================
\section{Background}
%======================================================================

\paragraph{Reformulations.}
$\max_{j\in\mathcal J} g_j(x)\le b \iff g_j(x)\le b\;\;\forall j\in\mathcal J$.

$\min_{x\in\mathcal X}\max_{j\in\mathcal J} g_j(x)=\min_{x\in\mathcal X,\,z}\, z$ s.t. $z\ge g_j(x)\;\;\forall j\in\mathcal J$.

\paragraph{Min-cost flow.}
$\min_{f}\ \sum_{(i,j)\in\mathcal E} c_{ij}f_{ij}$ s.t. $\sum_{j\in\delta^+(i)} f_{ij}-\sum_{j\in\delta^-(i)} f_{ji}=b_i\;\;\forall i\in\mathcal V$, $\;0\le f_{ij}\le u_{ij}\;\;\forall (i,j)\in\mathcal E$.

Theorem (integrality): If $b\in\mathbb Z^{\mathcal V}$ and $u\in\mathbb Z^{\mathcal E}$, then at least one optimal solution $f^*$ has integer coordinates.

Shortest $s$--$t$ path: $c_{ij}$ is the length of $(i,j)$, $b_s=1$, $b_t=-1$, $b_i=0$ otherwise, $u_{ij}=+\infty$.

Max $s$--$t$ flow: add edge $(t,s)$ with $c_{ts}=-1$, $u_{ts}=+\infty$, and set $c_{ij}=0\;\forall(i,j)\in\mathcal E$, $b_i=0\;\forall i\in\mathcal V$.

\paragraph{Bayes' rule.}
$q(y\mid x)=\dfrac{q(x,y)}{q(x)} = \dfrac{q(x\mid y)\,q(y)}{\sum_{\hat y\in Y} q(x\mid \hat y)\,q(\hat y)}$.
$\quad q(x) = \sum_{y\in Y}q(x\mid y)q(y)$. 
$\quad q(x,y) = q(x\mid y)q(y)$. 

\paragraph{Geometric series.}
$\sum_{t=0}^{+\infty}\delta^{t}=\dfrac{1}{1-\delta}$ and $\sum_{t=k}^{+\infty}\delta^{t}=\dfrac{\delta^{k}}{1-\delta}$ for $\delta\in[0,1)$.

%======================================================================
\section{Lookup}
%======================================================================

Action $a_i\in A_i$: one move.
Strategy $s_i$: a rule that gives an action for every type, history, or information set of player $i$, so $s_i(\theta_i)$ and $s_i(h)$ are actions.

\begin{center}
\begin{tabular}{@{}p{18mm}>{\raggedright\arraybackslash}p{74mm}>{\raggedright\arraybackslash}p{74mm}@{}}
\toprule
 & Complete information, moves observed & Private types or unobserved moves \\
\midrule
Static & Strategic form game: NE & Bayesian game: BNE \\
Dynamic & Extensive form game with perfect information, repeated game: SPE & Extensive form game with information sets: weak PBE \\
\bottomrule
\end{tabular}
\end{center}

\begin{center}
\footnotesize
\setlength{\tabcolsep}{4pt}
\renewcommand{\arraystretch}{1.4}
\begin{tabular}{@{}>{\raggedright\arraybackslash}p{25mm}>{\raggedright\arraybackslash}p{19mm}>{\raggedright\arraybackslash}p{76mm}>{\raggedright\arraybackslash}p{27mm}>{\raggedright\arraybackslash}p{27mm}@{}}
\toprule
Game & Equilibrium & Equilibrium payoff $\ge$ deviation payoff & Payoff conditioned on & For all \\
\midrule
Strategic form & NE $\sigma^*$ &
$u_i(\sigma_i^*,\sigma_{-i}^*)\ \ge\ u_i(\sigma_i,\sigma_{-i}^*)$ &
--- &
$\sigma_i\in\Sigma_i$, $i\in\mathcal I$ \\
Strategic form, correlated & CE $\pi^*$ &
$\sum_{s_{-i}\in S_{-i}}\pi^*(s_i,s_{-i})\,[u_i(s_i,s_{-i})-u_i(s_i',s_{-i})]\ \ge\ 0$ &
$s_i$, with $\pi^*(s_{-i}\mid s_i)$ &
$s_i,s_i'\in S_i$, $i\in\mathcal I$ \\
Bayesian & BNE $s^*$ &
$u_i(s_i^*(\theta_i),s_{-i}^*\mid\theta_i)\ \ge\ u_i(a_i,s_{-i}^*\mid\theta_i)$ &
$\theta_i$, with $p(\theta_{-i}\mid\theta_i)$ &
$a_i\in A_i$, $\theta_i\in\Theta_i$, $i\in\mathcal I$ \\
Extensive form, repeated & SPE $s^*$ &
$u_i(s_i^*\vert_h,s_{-i}^*\vert_h)\ \ge\ u_i(s_i,s_{-i}^*\vert_h)$ in $\Gamma(h)$ &
$h$ &
$s_i$ in $\Gamma(h)$, $h\in H$, $i\in\mathcal I$ \\
Extensive form, information sets & weak PBE $(s^*,\mu^*)$ &
$E_{\mu_i^*(\cdot\mid I)}[u_i(s_i^*\vert_I,s_{-i}^*\vert_I)]\ \ge\ E_{\mu_i^*(\cdot\mid I)}[u_i(s_i,s_{-i}^*\vert_I)]$ &
$I$, with $\mu_i^*(h\mid I)$ &
$s_i$, $I$ of player $i$, $i\in\mathcal I$ \\
Routing, nonatomic & WE $f$ &
$l_{p_1}(f)\ \le\ l_{p_2}(f)$, latency is a cost &
--- &
$p_1,p_2\in P_i$ with $f_{p_1}>0$, pair $i$ \\
Information design & IC signal $q$ &
$\sum_{\omega\in\Omega}q(a,\omega)\,[u(a,\omega)-u(a',\omega)]\ \ge\ 0$ &
recommendation $a$, with $\mu_a$ &
$a,a'\in A$ \\
\bottomrule
\end{tabular}
\end{center}

%======================================================================
\section{Strategic form game}
%======================================================================

\paragraph{Notation.}
$\Gamma = \langle\mathcal I, (S_i)_{i\in\mathcal I}, (u_i)_{i\in\mathcal I}\rangle$. $S=\prod_i S_i$. $s=(s_i,s_{-i})\in S$.
Strategy profile $\sigma=(\sigma_i)_{i\in\mathcal I}$: one distribution over $S_i$ per player. A pure profile $s$ is one action per player.

Mixed strategy $\sigma_i\in\Sigma_i = \{\sigma_i \in [0,1]^{S_i} : \sum_{s_i\in S_i}\sigma_i(s_i) = 1\}$.
Support $\operatorname{supp}(\sigma_i) = \{s_i\in S_i : \sigma_i(s_i) > 0\}$.
$u_i(\sigma) = E_{s\sim\sigma}[u_i(s)]$, with each $s_i\sim\sigma_i$ drawn independently.

Best response $BR_i(\sigma_{-i}) = \arg\max_{\sigma_i\in\Sigma_i} u_i(\sigma_i,\sigma_{-i})$.
$s_i$ is strictly dominated by $\hat s_i$ if $u_i(\hat s_i,s_{-i}) > u_i(s_i,s_{-i})\;\forall s_{-i}\in S_{-i}$.

Zero-sum game: $\mathcal I=\{1,2\}$ and $u_1(s_1,s_2) = -u_2(s_1,s_2)\;\forall(s_1,s_2)\in S_1\times S_2$.

Correlated strategy \new{$\pi\in\Delta(S)$}: a mediator draws $s\sim\pi$ and tells player $i$ only its recommendation $s_i$.

\paragraph{Properties.}
\begin{itemize}
    \item Equilibrium: No profitable unilateral \new{deviation}. Mutual consistency. Self-enforcement.
    \item \res{Lemma} Given $\sigma_{-i}$, at least one pure strategy is in $BR_i(\sigma_{-i})$. And $\sigma_i\in BR_i(\sigma_{-i}) \iff s_{i}\in BR_i(\sigma_{-i})\;\forall s_i\in\operatorname{supp}(\sigma_i)$.
    \item \res{Proposition} $\sigma^*$ is a Nash equilibrium \new{iff} $u_i(\sigma^*_i, \sigma^*_{-i}) \geq u_i(s_i,\sigma^*_{-i})\;\forall s_i\in S_i,\;\forall i\in\mathcal I$. Check pure deviations only.
    \item \res{Proposition} $\sigma^*$ is a Nash equilibrium \new{iff} $s_i\in BR_i(\sigma^*_{-i})\;\forall s_i\in\operatorname{supp}(\sigma^*_i),\;\forall i\in\mathcal I$. Then $u_i(s_i,\sigma^*_{-i})$ is equal across $s_i\in\operatorname{supp}(\sigma_i^*)$.
    \item \res{Theorem (Nash)} Every finite strategic-form game has at least one Nash equilibrium in mixed strategies.
    \item Strictly dominated strategies receive zero probability in every Nash Equilibrium. A pure strategy may be strictly dominated by a mixed strategy.
    \item Zero-sum. \res{Lemma} $(\sigma^*_1,\sigma^*_2)$ \new{is a Nash equilibrium iff} $\min_{\sigma_2\in\Sigma_2}u_1(\sigma^*_1,\sigma_2) = u_1(\sigma^*_1,\sigma^*_2) = \max_{\sigma_1\in\Sigma_1}\new{u_1}(\sigma_1,\sigma^*_2)$.
    \item Zero-sum. \res{Theorem (Von Neumann minmax)} $\max_{\sigma_1\in\Sigma_1}\min_{\sigma_2\in\Sigma_2}u_1(\sigma_1,\sigma_2) = \new{\min_{\sigma_2\in\Sigma_2}\max_{\sigma_1\in\Sigma_1}}u_1(\sigma_1,\sigma_2)$.
    \item Zero-sum. \res{Proposition} $(\sigma^*_1,\sigma^*_2)$ Nash Equilibrium $\iff \sigma^*_1\in\arg\max_{\sigma_1\in\Sigma_1}\min_{\sigma_2\in\Sigma_2}u_1(\sigma_1,\sigma_2)$ and $\sigma^*_2\in\new{\arg\min_{\sigma_2\in\Sigma_2}\max_{\sigma_1\in\Sigma_1}}u_1(\sigma_1,\sigma_2)$. The common optimal value $z^*=u_1(\sigma_1^*,\sigma_2^*)$ is the value of the game.
    \item Zero-sum. \res{Proposition} $(\sigma_1^*,z^*)$ and $(\sigma_2^*,z^*)$ optimally solve
    \begin{align*}
        \max_{\sigma_1\in\Sigma_1,\,z}\ z
        \;\;\text{s.t.}\;\; z\le\sum_{s_1\in S_1}\sigma_1(s_1)\,u_1(s_1,s_2)\;\;\forall s_2\in S_2
        \qquad\text{and}\qquad
        \min_{\sigma_2\in\Sigma_2,\,z}\ z
        \;\;\text{s.t.}\;\; z\ge\sum_{s_2\in S_2}\sigma_2(s_2)\,u_1(s_1,s_2)\;\;\forall s_1\in S_1.
    \end{align*}
    \item Correlated. Every mixed Nash equilibrium is a correlated equilibrium with $\pi(s)=\prod_{i\in\mathcal I}\sigma_i^*(s_i)$. A correlated equilibrium that is not a product is not a Nash equilibrium.
    \item Correlated. LP for a correlated equilibrium, with any linear objective:
    \begin{align*}
        \max_{\pi\in\Delta(S)}\ \sum_{s\in S}\pi(s)\sum_{i\in\mathcal I}u_i(s)
        \;\;\text{s.t.}\;\;
        \sum_{s_{-i}\in S_{-i}}\pi(s_i,s_{-i})\,[u_i(s_i,s_{-i})-u_i(s_i',s_{-i})]\ge 0\;\;\forall s_i,s_i'\in S_i,\;\forall i\in\mathcal I.
    \end{align*}
\end{itemize}

\paragraph{Problem format.}
Rows are player 1, columns are player 2, cells are $(u_1,u_2)$. Let $p=\sigma_1(C)$ and $q=\sigma_2(C)$.
\begin{center}
\begin{tabular}{c|cc}
 & $C$ & $A$ \\
\hline
$C$ & $6,6$ & $2,7$ \\
$A$ & $7,2$ & $0,0$ \\
\end{tabular}
\end{center}

\textbf{Strictly dominated actions.} An action is deleted if another action gives more against every action of the other player.
\begin{align*}
    u_1(A,C)=7>u_1(C,C)=6 \quad\text{but}\quad u_1(A,A)=0<u_1(C,A)=2.
\end{align*}
Neither row is strictly dominated. The same comparison holds for the columns.

\textbf{Pure NE.} $s_1^*\in BR_1(s_2^*)$ and $s_2^*\in BR_2(s_1^*)$.
\begin{align*}
    BR_1(C)=\{A\}\;\;(7>6),\qquad BR_1(A)=\{C\}\;\;(2>0),\qquad BR_2(C)=\{A\}\;\;(7>6),\qquad BR_2(A)=\{C\}\;\;(2>0).
\end{align*}
NE: $(C,A)$ and $(A,C)$.

\textbf{Mixed NE.} Each player is indifferent across its support.
\begin{align*}
    u_1(C,\sigma_2)=u_1(A,\sigma_2):&\quad 6q+2(1-q)=7q+0(1-q)\;\Rightarrow\;q=\tfrac23,\\
    u_2(\sigma_1,C)=u_2(\sigma_1,A):&\quad 6p+2(1-p)=7p+0(1-p)\;\Rightarrow\;p=\tfrac23,\\
    u_1(\sigma^*)=u_1(A,\sigma_2^*):&\quad 7q=\tfrac{14}{3}.
\end{align*}

\textbf{CE check} for $\pi(C,C)=\pi(C,A)=\pi(A,C)=\tfrac13$, $\pi(A,A)=0$, player 1.
\begin{align*}
    \text{After }C:&\quad \pi(C\mid C)=\frac{\pi(C,C)}{\pi(C,C)+\pi(C,A)}=\tfrac12,\qquad \tfrac12(6)+\tfrac12(2)=4\ \ge\ \tfrac12(7)+\tfrac12(0)=3.5,\\
    \text{After }A:&\quad \pi(C\mid A)=\frac{\pi(A,C)}{\pi(A,C)+\pi(A,A)}=1,\qquad 7\ \ge\ 6.
\end{align*}
Player 2 is symmetric, so $\pi$ is a CE.

%======================================================================
\section{Extensive form game}
%======================================================================

\paragraph{Notation.}
History $h^k=(a^1,\dots,a^{k-1})$: the actions played before stage $k$, with $h^1=\emptyset$.
A history is a node of the tree, named by the path that leads to it.
$H$ is the set of all histories.
Strategy profile $s=(s_i)_{i\in\mathcal I}$: one action per player at every nonterminal history where that player moves.

$A_i(h)$: the actions of player $i$ at $h$.
Pure strategy $s_i$: an action $s_i(h)\in A_i(h)$ at every nonterminal $h$, including $h$ off the outcome path.

Outcome of $s$: the terminal history that $s$ produces. $u_i(s)$ is player $i$'s payoff in the outcome of $s$.

Subgame $\Gamma(h)$: the game played after $h$, with payoffs inherited from $\Gamma$. $s\vert_h$ is the part of $s$ from $h$ onward.

Information set: histories that the moving player cannot distinguish. A subgame contains all of it or none of it.

Repeated game: stage game $\langle\mathcal I,(A_i)_{i\in\mathcal I},(g_i)_{i\in\mathcal I}\rangle$ played in periods $t=1,\dots,T$, with $T<+\infty$ in $\Gamma^T$ and $T=+\infty$ in $\Gamma^\infty$.
Players observe $h^t=(a^1,\dots,a^{t-1})$ before period $t$.
Strategy profile $s$: one action per player at every history $h^t$.
Pure strategy $s_i:H\to A_i$.
$u_i(s)=\sum_{t=1}^{T}\delta^{t-1}g_i(a^t(s))$ with $\delta\in[0,1)$, where $a^t(s)$ is the action profile that $s$ produces in period $t$.

\paragraph{Properties.}
\begin{itemize}
    \item Equivalent strategic form: the matrix whose rows and columns are the pure strategies $s_i$, with entries $u_i(s)$. A Nash equilibrium of $\Gamma$ is a Nash equilibrium of this matrix.
    \item SPE is a property of the full profile $s^*$: for every $h\in H$, $s^*\vert_h$ is a Nash equilibrium of $\Gamma(h)$.
    \item \res{One-stage deviation principle} For a finite-horizon game, or an infinite-horizon game with bounded discounted payoffs, $s^*$ is an SPE iff no player can gain by deviating from $s^*$ in a single stage and conforming to $s^*$ thereafter.
    \item \res{Theorem} Every finite horizon extensive form game with perfect information has a pure strategy SPE. Backward induction finds it.
    \item Repeated. \res{Theorem} If the stage game has a unique pure Nash equilibrium $a^*$, then $\Gamma^T$ has a unique SPE: $s_i^*(h^t)=a_i^*$ for every $h^t$ and every $i$.
    \item Repeated. Playing a stage-game Nash equilibrium after every history is an SPE of $\Gamma^T$ and of $\Gamma^\infty$.
    \item Repeated. Under a trigger strategy, player $i$'s stage payoff is $g_i^c$ in every period if it follows $s^*$. If it deviates once, its stage payoff is $g_i^d$ in that period and $g_i^p$ in every later period. Counting periods from the current one, player $i$ follows 
    $s^* 
    \iff\sum_{t=0}^{+\infty}\delta^t g_i^c\ \ge\ g_i^d+\sum_{t=1}^{+\infty}\delta^t g_i^p
    \quad\iff\quad
    \frac{g_i^c}{1-\delta}\ \ge\ g_i^d+\frac{\delta}{1-\delta}\,g_i^p.$
\end{itemize}

\paragraph{Problem format.}
Game tree. Player 1 chooses $A$ or $B$. After $A$, player 2 chooses $C\to(2,1)$ or $D\to(3,0)$. After $B$, player 2 chooses $E\to(4,2)$ or $F\to(1,3)$.

\textbf{Pure strategies.} One action at every history where the player moves.
\begin{align*}
    \text{Player 1 moves at }\emptyset:\ \{A,B\}.\qquad \text{Player 2 moves at }(A)\text{ and }(B):\ \{CE,\,CF,\,DE,\,DF\}.
\end{align*}

\textbf{Backward induction.} Start with the last subgames. Find the optimal action in each. Replace each solved subgame by its equilibrium payoffs. Continue toward the initial node.
\begin{align*}
    \Gamma(A):&\quad u_2(A,C)=1>u_2(A,D)=0\;\Rightarrow\;s_2^*((A))=C,\\
    \Gamma(B):&\quad u_2(B,F)=3>u_2(B,E)=2\;\Rightarrow\;s_2^*((B))=F,\\
    \Gamma(\emptyset):&\quad u_1(A,C)=2>u_1(B,F)=1\;\Rightarrow\;s_1^*(\emptyset)=A.
\end{align*}
SPE $s^*=(A,CF)$, outcome $(A,C)$, payoffs $(2,1)$.
A tie at a node gives one SPE per optimal action.

\textbf{Strategic form and pure NE.} Rows are player 1's pure strategies, columns are player 2's, and each cell is the payoff at the outcome.
\begin{center}
\begin{tabular}{c|cccc}
 & $CE$ & $CF$ & $DE$ & $DF$ \\
\hline
$A$ & $2,1$ & $2,1$ & $3,0$ & $3,0$ \\
$B$ & $4,2$ & $1,3$ & $4,2$ & $1,3$ \\
\end{tabular}
\end{center}
\begin{align*}
    &BR_1(CE)=\{B\},\quad BR_1(CF)=\{A\},\quad BR_1(DE)=\{B\},\quad BR_1(DF)=\{A\},\\
    &BR_2(A)=\{CE,CF\},\quad BR_2(B)=\{CF,DF\}.
\end{align*}
The only pure NE is $(A,CF)$, which is the SPE.

Repeated game. Stage game, rows are player 1.
\begin{center}
\begin{tabular}{c|cc}
 & Share & Withhold \\
\hline
Share & $1,1$ & $-1,2$ \\
Withhold & $2,-1$ & $0,0$ \\
\end{tabular}
\end{center}

\textbf{Stage game NE.} Withhold strictly dominates Share, $2>1$ and $0>-1$. NE: (Withhold, Withhold).

\textbf{$\Gamma^T$.} The stage game has a unique pure NE, so the unique SPE is $s_i^*(h^t)=$ Withhold for every $h^t$.
\begin{align*}
    u_i(s^*)=\sum_{t=1}^{T}\delta^{t-1}\cdot 0=0.
\end{align*}

\textbf{$\Gamma^\infty$, grim trigger.} $s_i^*(h^t)=$ Share if no player has withheld before period $t$, and Withhold otherwise.
Check one deviation in each case, counting periods from the current one.
\begin{align*}
    \text{No earlier Withhold}:&\quad \underbrace{\sum_{t=0}^{+\infty}\delta^t\cdot 1}_{\text{follow}}=\frac{1}{1-\delta}\ \ge\ \underbrace{2+\sum_{t=1}^{+\infty}\delta^t\cdot 0}_{\text{Withhold once}}=2\iff\delta\ge\tfrac12,\\
    \text{After a Withhold}:&\quad \underbrace{\sum_{t=0}^{+\infty}\delta^t\cdot 0}_{\text{follow}}=0\ \ge\ \underbrace{-1+\sum_{t=1}^{+\infty}\delta^t\cdot 0}_{\text{Share once}}=-1.
\end{align*}
Grim trigger is an SPE iff $\delta\ge\tfrac12$.

%======================================================================
\section{Bayesian game}
%======================================================================

\paragraph{Notation.}
$\Gamma=\langle\mathcal I,(A_i)_{i\in\mathcal I},(\Theta_i)_{i\in\mathcal I},p,(u_i)_{i\in\mathcal I}\rangle$.
Strategy profile $s=(s_i)_{i\in\mathcal I}$: one action per player per type.

Type $\theta_i\in\Theta_i$: what player $i$ knows and the other players do not.
$p$ is the prior over type profiles $\theta=(\theta_i,\theta_{-i})$.
$u_i(a\mid\theta)$ is player $i$'s payoff for action profile $a$ and type profile $\theta$.

Pure strategy $s_i:\Theta_i\to A_i$: an action for every type of player $i$, including the types that are not drawn.
Mixed strategy in $\Delta\big(A_i^{\Theta_i}\big)$: a distribution over pure strategies $s_i$.
Behavioral strategy $\sigma_i:\Theta_i\to\Delta(A_i)$: a distribution over $A_i$ for every type.

Interim payoff: the expected payoff of type $\theta_i$ when it plays the action $a_i$ and each other type $\theta_{-i}$ plays the action $s_{-i}(\theta_{-i})$.
$u_i(a_i,s_{-i}\mid\theta_i)=\sum_{\theta_{-i}\in\Theta_{-i}}p(\theta_{-i}\mid\theta_i)\cdot u_i\big(a_i,s_{-i}(\theta_{-i})\mid\theta_i,\theta_{-i}\big)$.

\paragraph{Properties.}
\begin{itemize}
    \item Timing, which is the equivalent extensive form: Nature draws $\theta\sim p$. Each player $i$ learns $\theta_i$ only. All players then choose $a_i$ at the same time. Payoffs are $u_i(a\mid\theta)$.
    \item BNE is a property of the full profile $s^*$: for every $i\in\mathcal I$ and every $\theta_i\in\Theta_i$, $s_i^*(\theta_i)\in\arg\max_{a_i\in A_i}u_i(a_i,s_{-i}^*\mid\theta_i)$.
    \item \res{Assumption} The action spaces, payoff functions, possible types, and the prior are common knowledge.
    \item \res{Theorem} Every Bayesian game with finitely many players, types, and actions has a Bayesian Nash equilibrium in mixed strategies.
    \item \res{Proposition (Kuhn)} With perfect recall, every mixed strategy has an outcome-equivalent behavioral strategy, and conversely.
    \item With mixing: $\operatorname{supp}\sigma_i^*(\cdot\mid\theta_i) \subseteq \arg\max_{a_i\in A_i}u_i(a_i,\sigma_{-i}^*\mid\theta_i)\;\forall\theta_i\in\Theta_i,\;\forall i\in\mathcal I$. A type mixes only if it is indifferent.
    \item Information can lower equilibrium payoffs by changing others' incentives.
\end{itemize}

\paragraph{Problem format.}
Organization 1 is type $T$ or $F$, each with probability $\tfrac12$. Organization 2 has one type. Rows are $a_1$, columns are $a_2$, cells are $(u_1,u_2)$.
\begin{center}
\begin{tabular}{c|cc}
\multicolumn{3}{c}{Type $T$}\\
 & $R$ & $N$ \\
\hline
$R$ & $16,18$ & $6,14$ \\
$N$ & $14,8$ & $0,0$ \\
\end{tabular}
\qquad\qquad
\begin{tabular}{c|cc}
\multicolumn{3}{c}{Type $F$}\\
 & $R$ & $N$ \\
\hline
$R$ & $-8,-6$ & $-8,0$ \\
$N$ & $0,-6$ & $0,0$ \\
\end{tabular}
\end{center}

\textbf{Pure strategies.} One action per type.
\begin{align*}
    s_1=(s_1(T),s_1(F))\in\{RR,\,RN,\,NR,\,NN\},\qquad s_2\in\{R,N\},\qquad 4\times 2=8\text{ pure profiles}.
\end{align*}

\textbf{Beliefs.} Organization 1 knows $\theta_1$. Organization 2 has one type, so it uses the prior, $p(T)=p(F)=\tfrac12$.

\textbf{Each type of organization 1.} $s_1^*(\theta_1)\in\arg\max_{a_1}u_1(a_1,s_2^*\mid\theta_1)$.
\begin{align*}
    \text{Type }T:&\quad u_1(R,R\mid T)=16>u_1(N,R\mid T)=14,\qquad u_1(R,N\mid T)=6>u_1(N,N\mid T)=0\;\Rightarrow\;s_1^*(T)=R,\\
    \text{Type }F:&\quad u_1(N,R\mid F)=0>u_1(R,R\mid F)=-8,\qquad u_1(N,N\mid F)=0>u_1(R,N\mid F)=-8\;\Rightarrow\;s_1^*(F)=N.
\end{align*}
Each type has a strictly dominant action, so these hold for either $s_2$.

\textbf{Organization 2.} $s_2^*\in\arg\max_{a_2}\sum_{\theta_1}p(\theta_1)\,u_2\big(s_1^*(\theta_1),a_2\mid\theta_1\big)$.
\begin{align*}
    a_2=R:&\quad \tfrac12\,u_2(R,R\mid T)+\tfrac12\,u_2(N,R\mid F)=\tfrac12(18)+\tfrac12(-6)=6,\\
    a_2=N:&\quad \tfrac12\,u_2(R,N\mid T)+\tfrac12\,u_2(N,N\mid F)=\tfrac12(14)+\tfrac12(0)=7\;\Rightarrow\;s_2^*=N.
\end{align*}
BNE: $s_1^*(T)=R$, $s_1^*(F)=N$, $s_2^*=N$.

\textbf{Mixed BNE.} As for a mixed NE, type by type. Each type $\theta_i$ that mixes is indifferent: $u_i(a_i,\sigma_{-i}\mid\theta_i)=u_i(a_i',\sigma_{-i}\mid\theta_i)\;\forall a_i,a_i'\in\operatorname{supp}\sigma_i(\cdot\mid\theta_i)$. Solve these equations for $\sigma_{-i}$.


%======================================================================
\section{Extensive form game with information sets}
%======================================================================

\paragraph{Notation.}
Information set $I$: a set of histories that the player who moves there cannot tell apart. That player chooses one action for all of $I$.
Strategy profile $s$: one action per player at every information set of that player.

Belief $\mu_i(h\mid I)$: the probability that player $i$ assigns to being at history $h$, given that $I$ is reached.
Assessment $(s,\mu)$: a strategy profile $s$ and a belief $\mu_i(\cdot\mid I)$ at every information set.

On-path $I$: $s$ reaches $I$ with positive probability. Off-path $I$: $s$ reaches $I$ with probability $0$.

$s\vert_I$ is the part of $s$ from $I$ onward.
$E_{\mu_i(\cdot\mid I)}[u_i(s\vert_I)]$ is player $i$'s expected payoff at $I$: draw $h\in I$ from $\mu_i(\cdot\mid I)$, then follow $s$ from $h$.

Signaling game, timing: Nature draws the sender's type $\theta_1\sim p$. The sender (player 1) learns $\theta_1$ and sends a message $m$. The receiver (player 2) sees $m$, not $\theta_1$, and chooses an action $a$. Payoffs are $u_1(m,a\mid\theta_1)$ and $u_2(m,a\mid\theta_1)$.
Strategy profile $s=(s_1,s_2)$: $s_1(\theta_1)$ is the message of type $\theta_1$, and $s_2(m)$ is the receiver's action after $m$.
The sender has one information set per type. The receiver has one per message, with belief $\mu_2(\theta_1\mid m)$.

\paragraph{Properties.}
\begin{itemize}
    \item Weak PBE is a property of an assessment $(s^*,\mu^*)$. It has two parts, sequential rationality and belief consistency.
    \item Sequential rationality: at every $I$, on-path or off-path, the player $i$ who moves at $I$ plays a best response given its belief $\mu_i^*(\cdot\mid I)$ and given that all players follow $s^*$ after $I$.
    $E_{\mu_i^*(\cdot\mid I)}[u_i(s_i^*\vert_I,s_{-i}^*\vert_I)]\ \ge\ E_{\mu_i^*(\cdot\mid I)}[u_i(s_i,s_{-i}^*\vert_I)]\;\;\forall s_i$.
    \item Belief consistency: at every on-path $I$, $\mu_i^*(\cdot\mid I)$ comes from $p$ and $s^*$ by Bayes' rule. At an off-path $I$, Bayes' rule divides by $0$, so every belief is allowed.
    \item SPE tests $s^*$ only at a history where a subgame starts, and no subgame starts inside an $I$ with two or more histories. Weak PBE tests $s^*$ at every $I$.
    \item Signaling. A weak PBE is $(s_1^*,s_2^*,\mu_2^*)$ with three conditions.
    \begin{align*}
        \text{Receiver, after every } m:\quad & s_2^*(m)\in\arg\max_{a}\sum_{\theta_1\in\Theta_1}\mu_2^*(\theta_1\mid m)\,u_2(m,a\mid\theta_1).\\
        \text{Sender, for every } \theta_1:\quad & s_1^*(\theta_1)\in\arg\max_{m}u_1\big(m,s_2^*(m)\mid\theta_1\big).\\
        \text{Belief, after every sent } m:\quad & \mu_2^*(\theta_1\mid m)=\frac{P(m\mid\theta_1)\,p(\theta_1)}{\sum_{\theta_1'\in\Theta_1}P(m\mid\theta_1')\,p(\theta_1')},\quad P(m\mid\theta_1)=1\text{ if }s_1^*(\theta_1)=m\text{ and }0\text{ otherwise}.
    \end{align*}
    \item Separating: each type sends a different $m$, so $\mu_2^*(\theta_1\mid m)\in\{0,1\}$ after every sent $m$. Every $m$ is on-path.
    \item Pooling: both types send the same $m$, so $\mu_2^*(\theta_1\mid m)=p(\theta_1)$. The unsent message is off-path.
    \item A signal is credible when producing it is sufficiently more costly for the type that would benefit from the deception.
\end{itemize}

\paragraph{Problem format.}
Sender type $H$ or $L$ with $p(H)=\tfrac13$. Rows are the message $m$, columns are the receiver's action $a$, cells are $(u_1,u_2)$.
\begin{center}
\begin{tabular}{c|cc}
\multicolumn{3}{c}{Type $H$}\\
 & $D$ & $R$ \\
\hline
$C$ & $2,2$ & $-1,0$ \\
$N$ & $3,2$ & $0,0$ \\
\end{tabular}
\qquad\qquad
\begin{tabular}{c|cc}
\multicolumn{3}{c}{Type $L$}\\
 & $D$ & $R$ \\
\hline
$C$ & $-1,-2$ & $-4,0$ \\
$N$ & $3,-2$ & $0,0$ \\
\end{tabular}
\end{center}

\textbf{Pooling on $N$}, $s_1(H)=s_1(L)=N$.
\begin{align*}
    \text{Belief after }N:&\quad \mu_2(H\mid N)=\frac{P(N\mid H)\,p(H)}{P(N\mid H)\,p(H)+P(N\mid L)\,p(L)}=\frac{1\cdot\frac13}{1\cdot\frac13+1\cdot\frac23}=\tfrac13,\\
    \text{Receiver after }N:&\quad \underbrace{\tfrac13(2)+\tfrac23(-2)}_{D}=-\tfrac23\ <\ \underbrace{0}_{R}\;\Rightarrow\;s_2(N)=R,\\
    \text{Receiver after }C\text{, off-path}:&\quad \text{with }D\text{ after }C\text{, type }H\text{ gets }u_1(C,D\mid H)=2>u_1(N,R\mid H)=0\text{, so }s_2(C)=R\text{ is needed},\\
    \text{Belief after }C:&\quad R\text{ is a best response iff }\underbrace{2\mu_2(H\mid C)-2\big(1-\mu_2(H\mid C)\big)}_{D}\le\underbrace{0}_{R}\iff\mu_2(H\mid C)\le\tfrac12,\\
    \text{Type }H:&\quad u_1(N,R\mid H)=0\ \ge\ u_1(C,R\mid H)=-1,\\
    \text{Type }L:&\quad u_1(N,R\mid L)=0\ \ge\ u_1(C,R\mid L)=-4.
\end{align*}
Weak PBE: $s_1^*(H)=s_1^*(L)=N$, $s_2^*(C)=s_2^*(N)=R$, $\mu_2^*(H\mid N)=\tfrac13$, $\mu_2^*(H\mid C)\le\tfrac12$.

\textbf{Separating}, $s_1(H)=C$, $s_1(L)=N$.
\begin{align*}
    \text{Beliefs}:&\quad \mu_2(H\mid C)=\frac{1\cdot\frac13}{1\cdot\frac13+0\cdot\frac23}=1,\qquad \mu_2(H\mid N)=\frac{0\cdot\frac13}{0\cdot\frac13+1\cdot\frac23}=0,\\
    \text{Receiver after }C:&\quad u_2(C,D\mid H)=2>u_2(C,R\mid H)=0\;\Rightarrow\;s_2(C)=D,\\
    \text{Receiver after }N:&\quad u_2(N,D\mid L)=-2<u_2(N,R\mid L)=0\;\Rightarrow\;s_2(N)=R,\\
    \text{Type }H:&\quad u_1(C,D\mid H)=2\ \ge\ u_1(N,R\mid H)=0,\\
    \text{Type }L:&\quad u_1(N,R\mid L)=0\ \ge\ u_1(C,D\mid L)=-1.
\end{align*}
Weak PBE: $s_1^*(H)=C$, $s_1^*(L)=N$, $s_2^*(C)=D$, $s_2^*(N)=R$, $\mu_2^*(H\mid C)=1$, $\mu_2^*(H\mid N)=0$.

\textbf{Semi-separating}. One type sends one message, and the other type mixes with probability $x$. Receiver indifference after $m$ determines $x$. The receiver then mixes with probability $y$, and sender indifference determines $y$. Check the type that does not mix.


%======================================================================
\section{Routing game}
%======================================================================
\todo{NEED TO REVIEW}
\paragraph{Notation.}
Traffic instance $(G,r,l)$: directed graph $G=(\mathcal V,\mathcal E)$, rate $r_i>0$ for each origin-destination pair $i$, and a nonnegative, nondecreasing latency $l_e(\cdot)$ on each edge.
Players are nonatomic: each driver is too small to change the flow.

Flow $f$: the mass $f_p$ on each path $p\in P_i$ of each pair $i$, with $\sum_{p\in P_i}f_p=r_i$. It is the strategy profile.
Edge flow $f_e=\sum_{\{p\,:\,e\in p\}}f_p$.

Path latency $l_p(f)=\sum_{e\in p}l_e(f_e)$.
Cost $C(f)=\sum_{e\in\mathcal E}l_e(f_e)\,f_e$.
Marginal cost $l_e^*(f_e)=\frac{d}{df_e}\big[l_e(f_e)\,f_e\big]=l_e(f_e)+l_e'(f_e)\,f_e$, and $l_p^*(f)=\sum_{e\in p}l_e^*(f_e)$.

\paragraph{Properties.}
\begin{itemize}
    \item Wardrop equilibrium (WE) is a property of a feasible flow $f$: $l_{p_1}(f)\le l_{p_2}(f)$ for every pair $i$ and every $p_1,p_2\in P_i$ with $f_{p_1}>0$. The used paths of a pair have equal latency. It is the Nash equilibrium of the nonatomic game.
    \item \res{Lemma} A feasible $f$ is socially optimal, meaning it minimizes $C(f)$, iff $l_{p_1}^*(f)\le l_{p_2}^*(f)$ for every pair $i$ and every $p_1,p_2\in P_i$ with $f_{p_1}>0$.
    \item \res{Corollary} If $l_e(f_e)\,f_e$ is convex for every $e$, $f$ is socially optimal for $(G,r,l)$ iff it is a WE for $(G,r,l^*)$.
    \item $l_e'(f_e)\,f_e$ is the externality that one more unit of flow puts on the flow already on $e$. A toll equal to it on each edge makes the WE socially optimal.
    \item With $l_e$ continuous, $f$ is a WE iff it minimizes $\sum_{e\in\mathcal E}\int_0^{f_e}l_e(z)\,dz$ over feasible flows. A WE exists, and all WE have the same $C(f)$.
    \item Braess paradox: a new edge never increases the socially optimal cost, and it can increase the WE cost.
\end{itemize}

\paragraph{Problem format.}
Pigou network. A unit mass chooses path 1 with $l_1(f_1)=f_1$ or path 2 with $l_2=1$.
\begin{align*}
    \textbf{WE}:&\quad l_1(f_1)=f_1<1=l_2\text{ for every }f_1<1\;\Rightarrow\;f_1=1,\qquad C=1\cdot 1=1,\\
    \textbf{Social optimum}:&\quad l_1^*(f_1)=2f_1=1=l_2^*\;\Rightarrow\;f_1=\tfrac12,\qquad C=\tfrac12\cdot\tfrac12+\tfrac12\cdot 1=\tfrac34.
\end{align*}

%======================================================================
\section{Information design}
%======================================================================
\todo{NEED TO REVIEW}
\paragraph{Notation.}
Sender and Receiver. State $\omega\in\Omega$ with common prior $\mu_0$. Receiver chooses $a\in A$. Payoffs are $u(a,\omega)$ for Receiver and $v(a,\omega)$ for Sender.

Signal $\pi$: a distribution $\pi(\cdot\mid\omega)$ over realizations $s\in S$ for every state.
Timing: Sender publicly commits to $\pi$. Nature draws $\omega\sim\mu_0$ and $s\sim\pi(\cdot\mid\omega)$. Receiver sees $s$, not $\omega$, and chooses $a$.

Posterior $\mu_s(\omega)=\dfrac{\mu_0(\omega)\,\pi(s\mid\omega)}{P(s)}$, with $P(s)=\sum_{\omega\in\Omega}\mu_0(\omega)\,\pi(s\mid\omega)$.

$\hat a(\mu)\in\arg\max_{a\in A}E_\mu[u(a,\omega)]$: Receiver's action at belief $\mu$, with ties broken in Sender's favor.
$\hat v(\mu)=E_\mu\big[v(\hat a(\mu),\omega)\big]$: Sender's payoff at belief $\mu$.

$\tau$: the distribution of posteriors that $\pi$ produces. It puts probability $P(s)$ on $\mu_s$.
$\tau$ is Bayes plausible if $E_\tau[\mu]=\mu_0$.
$V$: the concave closure of $\hat v$, the upper boundary of the convex hull of its graph.

Straightforward signal: each realization is a recommended action, and Receiver's best action is to follow it. $q(a,\omega)=\mu_0(\omega)\,\pi(a\mid\omega)$.

Cheap talk: no commitment. Sender chooses $s$ after it observes $\omega$.

Many receivers, Bayesian routing game: state $\theta\in\{0,1\}$ with $P(\theta=1)=p$. Each driver gets a private recommendation $r\in\{1,2\}$, drawn independently given $\theta$, with $q_\theta=P(r=2\mid\theta)$. If all drivers follow, the flow on path 2 in state $\theta$ is $q_\theta$.

\paragraph{Properties.}
\begin{itemize}
    \item Straightforward signals are without loss of generality. For finite $A$, an optimal signal solves the LP below. Its second constraint is incentive compatibility (IC).
    \begin{align*}
        \max_{q\ge 0}\ \sum_{a\in A}\sum_{\omega\in\Omega}q(a,\omega)\,v(a,\omega)
        \;\;\text{s.t.}\;\;
        \sum_{a\in A}q(a,\omega)=\mu_0(\omega)\;\;\forall\omega\in\Omega,
        \qquad
        \sum_{\omega\in\Omega}q(a,\omega)\,[u(a,\omega)-u(a',\omega)]\ge 0\;\;\forall a,a'\in A.
    \end{align*}
    \item Every signal produces a Bayes-plausible $\tau$. Every finitely supported Bayes-plausible $\tau$ is produced by the signal $\pi(s\mid\omega)=\dfrac{\mu_s(\omega)\,\tau(\mu_s)}{\mu_0(\omega)}$.
    \item \res{Proposition (Kamenica and Gentzkow)} Sender's problem is $\max_\tau E_\tau[\hat v(\mu)]$ s.t. $E_\tau[\mu]=\mu_0$. \res{Proposition} Its value is $V(\mu_0)$.
    \item Sender benefits from persuasion iff $V(\mu_0)>\hat v(\mu_0)$. If $\hat v$ is concave, Sender does not benefit. If $\hat v$ is convex, full revelation is optimal.
    \item Any signal weakly improves Receiver's expected payoff relative to no information.
    \item Cheap-talk equilibrium: a signal $\pi$ with Bayes posteriors $\mu_s$, Receiver playing $\hat a(\mu_s)$, and Sender IC. Sender IC is $s\in\arg\max_{s'\in S}v\big(\hat a(\mu_{s'}),\omega\big)$ for every $\omega\in\Omega$ and every $s\in\operatorname{supp}\pi(\cdot\mid\omega)$.
    \item Many receivers. A recommendation rule is a Bayes correlated equilibrium if following is optimal for a driver after every recommendation, given its posterior and given that the other drivers follow.
    \item Bayesian routing example. \res{Proposition} The socially optimal rule $(q_0,q_1)=(0,\tfrac16)$ satisfies the IC after $r=2$ for every $p$, and the IC after $r=1$ iff $p\le\tfrac{6}{11}$.
\end{itemize}

\paragraph{Problem format.}
School and company. $\Omega=\{G,M\}$, $\mu_0(G)=\tfrac12$, and $\mu=\mu(G)$. Hiring ($a=1$) gives the company $1$ if $G$ and $-2$ if $M$, and gives the school $1$. Rejecting ($a=0$) gives both $0$.
The exam is the signal $\pi(+\mid G)=1$, $\pi(+\mid M)=\tfrac12$.
\begin{align*}
    \textbf{Threshold}:&\quad \mu\cdot 1+(1-\mu)(-2)=3\mu-2\ge 0\iff\mu\ge\tfrac23,\qquad \hat v(\mu)=1\text{ if }\mu\ge\tfrac23,\ 0\text{ otherwise},\\
    \textbf{Value of the exam}:&\quad P(+)=\tfrac12\cdot 1+\tfrac12\cdot\tfrac12=\tfrac34,\qquad \mu_+=\frac{\tfrac12\cdot 1}{\tfrac34}=\tfrac23,\qquad \mu_-=0,\qquad E_\tau[\hat v(\mu)]=\tfrac34\cdot 1+\tfrac14\cdot 0=\tfrac34,\\
    \textbf{Concavification}:&\quad V(\mu)=\min\{\tfrac32\mu,1\},\qquad V(\mu_0)=\tfrac34,\qquad \text{endpoints }\mu_-=0,\ \mu_+=\tfrac23,\\
    \textbf{Bayes plausibility}:&\quad \tau(\mu_+)\cdot\tfrac23+\tau(\mu_-)\cdot 0=\tfrac12\;\Rightarrow\;\tau(\mu_+)=\tfrac34,\\
    \textbf{Signal from }\tau:&\quad \pi(+\mid G)=\frac{\mu_+(G)\,\tau(\mu_+)}{\mu_0(G)}=\frac{\tfrac23\cdot\tfrac34}{\tfrac12}=1,\qquad \pi(+\mid M)=\frac{\tfrac13\cdot\tfrac34}{\tfrac12}=\tfrac12,\\
    \textbf{LP}:&\quad \max_{q\ge 0}\ q(1,G)+q(1,M)\;\;\text{s.t.}\;\;q(1,G)+q(0,G)=\tfrac12,\quad q(1,M)+q(0,M)=\tfrac12,\\
    &\quad q(1,G)\cdot 1+q(1,M)\cdot(-2)\ge 0,\qquad q(0,G)\cdot(-1)+q(0,M)\cdot 2\ge 0,\\
    \textbf{Cheap talk}:&\quad -\in\operatorname{supp}\pi(\cdot\mid M)\text{ and }v\big(\hat a(\mu_-),M\big)=0<1=v\big(\hat a(\mu_+),M\big)\text{, so Sender IC fails}.
\end{align*}
School's payoff: $0$ with no information, $\tfrac12$ with full revelation, $\tfrac34$ with the optimal signal. Company's payoff: $0$, $\tfrac12$, $0$.

Bayesian routing game. $l_1=\theta$ and $l_2(x)=\tfrac13+2x$, with $x$ the flow on path 2. Cost $C(x,\theta)=(1-x)\,\theta+x\,\big(\tfrac13+2x\big)$. Take $p=\tfrac12$.
\begin{align*}
    \textbf{Social optimum}:&\quad \theta=0:\ x=0.\qquad \theta=1:\ \tfrac{d}{dx}C(x,1)=-\tfrac23+4x=0\;\Rightarrow\;x=\tfrac16,\qquad E_\theta[C]=\tfrac12\cdot\tfrac{17}{18}=\tfrac{17}{36},\\
    \textbf{Full information}:&\quad \theta=0:\ x=0.\qquad \theta=1:\ \tfrac13+2x=1\;\Rightarrow\;x=\tfrac13,\qquad E_\theta[C]=\tfrac12\cdot 1=\tfrac12,\\
    \textbf{No information}:&\quad \tfrac13+2x=E[\theta]=\tfrac12\;\Rightarrow\;x=\tfrac{1}{12},\qquad E_\theta[C]=\tfrac12,\\
    \textbf{IC after }r=1:&\quad (1-p)(1-q_0)\big(0-\tfrac13-2q_0\big)+p\,(1-q_1)\big(1-\tfrac13-2q_1\big)\le 0,\\
    \textbf{IC after }r=2:&\quad (1-p)\,q_0\big(\tfrac13+2q_0-0\big)+p\,q_1\big(\tfrac13+2q_1-1\big)\le 0,\\
    \textbf{Check }(q_0,q_1)=(0,\tfrac16):&\quad \tfrac12\cdot 1\cdot\big(-\tfrac13\big)+\tfrac12\cdot\tfrac56\cdot\tfrac13=-\tfrac{1}{36}\le 0,\qquad 0+\tfrac12\cdot\tfrac16\cdot\big(-\tfrac13\big)=-\tfrac{1}{36}\le 0,\\
    \textbf{Platform}:&\quad \min_{(q_0,q_1)\in[0,1]^2}\ (1-p)\,q_0\big(\tfrac13+2q_0\big)+p\,\big[(1-q_1)+q_1\big(\tfrac13+2q_1\big)\big]\;\;\text{s.t. both IC}.
\end{align*}
Each IC is the driver's expected latency on the recommended path minus that on the other path, multiplied by $P(r)$.


\section{Remaining}
\todo{Learning, learning from actions}
\todo{Planning with a model}

\end{document}